% =====================================================================
% paper.tex  —  From Natural Language to Verifiable Reasoning:
%               A Semantic Reduction Framework for Reliable AI Mathematics
%
% Submission-ready LaTeX source (pdflatex / xelatex compatible).
% Compile:  pdflatex paper && bibtex paper && pdflatex paper && pdflatex paper
%
% All factual claims in this manuscript are grounded in the cited literature;
% AI-assisted drafting is logged in AI日志.md.
% =====================================================================

\documentclass[11pt]{article}

\usepackage[utf8]{inputenc}
\usepackage{amsmath}
\usepackage{amssymb}
\usepackage{booktabs}
\usepackage{url}
\usepackage{hyperref}
\hypersetup{colorlinks=true, citecolor=blue, linkcolor=blue, urlcolor=blue}

\title{From Natural Language to Verifiable Reasoning:\\
A Semantic Reduction Framework for Reliable AI Mathematics}

\author{贾静文\\
郑州西亚斯学院\\
\texttt{3881471857@qq.com}}

\date{}

\begin{document}

\maketitle

% ---------------------------------------------------------------------
\begin{abstract}
Large language models (LLMs) can now produce mathematical reasoning that \emph{looks}
correct yet is frequently wrong. Reliability in AI for Mathematics (AI4Math) therefore
cannot rest on the plausibility of natural-language (NL) text alone. We argue that
trustworthy mathematical AI requires an explicit \emph{semantic reduction layer} that
translates NL reasoning into a typed, machine-checkable intermediate representation
(IR) before verification. We formalize a three-layer model---Generation, Semantic
Reduction, Verification---and identify the semantic reduction layer as the central
bottleneck: most failures originate there, not in generation or verification. We then
propose a concrete framework built on a \emph{Typed IR} (with a type system, operational
semantics, constraints, and a compositional structure) and a four-stage pipeline
(NL $\rightarrow$ IR $\rightarrow$ verification $\rightarrow$ NL rendering) with a
bidirectional-consistency requirement that guarantees interpretability and
traceability. Beyond standard accuracy, we define two additional reliability metrics---
\emph{verifiability} and \emph{consistency}---that bound what a correctness claim may
legitimately assert. We map our framework onto real systems (Draft--Sketch--Prove,
LeanDojo/ReProver, Math-Shepherd, AlphaGeometry) and distill a practical
\emph{reliability checklist}. This is a framework / position paper: the contribution is
conceptual and methodological, accompanied by a validation protocol rather than a new
large-scale empirical result.
\end{abstract}

% ---------------------------------------------------------------------
\section{Introduction}
\label{sec:intro}

\paragraph{Background.}
Large language models such as GPT-4 and Claude have demonstrated surprising competence
on mathematical reasoning benchmarks when paired with chain-of-thought (CoT)
prompting~\cite{wei2022chain} and search/self-consistency
strategies~\cite{yao2023tree}. Yet the same models routinely emit proofs and derivations
that are fluent, locally plausible, and globally invalid---a phenomenon often called
\emph{hallucination} in mathematical contexts.

\paragraph{Problem.}
Natural-language reasoning is \emph{not directly verifiable}. A paragraph of CoT has no
formal semantics: nothing in the text guarantees that each inference step preserves
truth, that defined terms are used consistently, or that the final answer follows from
the premises. This is the core \emph{reliability} (or trustworthiness) problem of
AI4Math.

\paragraph{Current approaches.}
Today's systems address reliability in two ways. The first is \emph{pure generation}:
train or prompt an LLM to reason and hope the answer is right. The second is
\emph{post-hoc verification}: check the output after the fact using proof assistants
(Coq, Lean, Agda), symbolic solvers (SMT), or learned verifiers / reward
models~\cite{cobbe2021training,lightman2024lets,wang2024mathshepherd}.

\paragraph{Thesis.}
We contend that reliability cannot rely \emph{solely} on final verification. Post-hoc
checking only works if the claim has already been expressed in a form a checker can
consume; if the gap between NL and that form is unstable, verification has no anchor.
We therefore argue for an intermediate \emph{semantic reduction layer} as a first-class
system component.

\paragraph{Contributions.}
\begin{enumerate}
    \item We formalize a \textbf{three-layer model} of reliable AI mathematics
          (Generation / Semantic Reduction / Verification) and argue, with a concrete
          failure analysis, that the \textbf{semantic reduction layer is the central
          bottleneck}.
    \item We propose a \textbf{Typed IR} design (type system, operational semantics,
          constraints, compositional structure) and a four-stage
          \textbf{pipeline} with a \textbf{bidirectional-consistency} requirement.
    \item We define two reliability metrics beyond accuracy---
          \textbf{verifiability} and \textbf{consistency}---that make the boundary of
          any correctness claim explicit.
    \item We synthesize a practical \textbf{reliability checklist} mapping each threat
          class (search, formalization, evaluation) to a concrete mechanism.
\end{enumerate}

% ---------------------------------------------------------------------
\section{A Three-Layer Model of Reliable AI Mathematics}
\label{sec:model}

We decompose a reliable AI mathematics system into three layers. The structure is
illustrated in Figure~\ref{fig:layers}.

\begin{figure}[t]
\centering
\fbox{\parbox{0.92\linewidth}{\centering
\textbf{Generation Layer} \\ (LLM, Natural-Language Chain-of-Thought) \\[2pt]
\small probabilistic $\cdot$ untyped $\cdot$ informal
}}
\\[5pt] $\downarrow$ \\[5pt]
\fbox{\parbox{0.92\linewidth}{\centering
\textbf{Semantic Reduction Layer} \\ (Typed IR: DSL / AST / Proof Object / Workflow Graph) \\[2pt]
\small \emph{--- bottleneck ---} ambiguity $\cdot$ missing typing $\cdot$ non-invertibility
}}
\\[5pt] $\downarrow$ \\[5pt]
\fbox{\parbox{0.92\linewidth}{\centering
\textbf{Verification Layer} \\ (Coq / Lean / Agda / SMT / symbolic solvers) \\[2pt]
\small consistency $\cdot$ derivability $\cdot$ factual grounding $\cdot$ constraint satisfaction
}}
\caption{The three-layer model. Reliability fails most often at the \emph{semantic
reduction layer} (middle), not at generation or verification.}
\label{fig:layers}
\end{figure}

\subsection{Generation Layer}
\label{sec:gen}
The generation layer produces natural-language reasoning---typically a chain of
thought~\cite{wei2022chain} or a tree of thought~\cite{yao2023tree}. Its defining
properties are that it is \textbf{probabilistic} (outputs are sampled, not derived),
\textbf{untyped} (no formal guarantee that objects are well-formed), and
\textbf{informal} (no binding semantics). It is the interface through which humans
state problems and read answers, but by itself it provides no reliability guarantee.

\subsection{Semantic Reduction Layer}
\label{sec:reduction}
The semantic reduction layer translates natural language into a structured, typed
representation. Candidate target forms include:
\begin{itemize}
    \item \textbf{Logical formulae} (first-order / higher-order logic);
    \item \textbf{Programs} (Python or a domain-specific language, DSL);
    \item \textbf{Proof objects} (e.g., Lean/Coq proof terms);
    \item \textbf{Typed intermediate representations (IR)} / abstract syntax trees
          (AST) / workflow graphs.
\end{itemize}
Its theoretical grounding is the \textbf{Curry--Howard correspondence}
\cite{howard1980formulae,sorensen2006lectures}, which identifies proofs with typed
programs and propositions with types---precisely the bridge that lets an informal
argument become a checkable object. This layer, however, faces three intrinsic
challenges:
\begin{itemize}
    \item \textbf{Ambiguity}: a single NL sentence may admit several formal readings;
    \item \textbf{Missing typing}: NL rarely states the types of objects, so the IR
          must recover or impose them;
    \item \textbf{Non-invertibility}: mapping NL $\rightarrow$ IR is lossy and not
          uniquely reversible, so traceability must be engineered explicitly.
\end{itemize}

\subsection{Verification Layer}
\label{sec:verify}
The verification layer checks the reduced object in a formal system. It establishes:
\begin{itemize}
    \item \textbf{Consistency} (no contradiction is introduced);
    \item \textbf{Derivability} (the conclusion follows from the premises under the
          system's inference rules);
    \item \textbf{Factual grounding} (constants and references match an accepted body
          of mathematics);
    \item \textbf{Constraint satisfaction} (stated side-conditions hold).
\end{itemize}
Tools range from proof assistants (Coq, Agda, Lean) to SMT and symbolic solvers.

% ---------------------------------------------------------------------
\section{The Bottleneck: Why Semantic Reduction Fails First}
\label{sec:bottleneck}

We now make the central empirical claim of this paper: \textbf{most reliability
failures in current AI4Math systems originate in the semantic reduction layer, not in
generation or verification.}

\paragraph{Argument.}
Verification depends on a \emph{computable object} it can manipulate. Generation can be
made more fluent and verification can be made stricter, but neither helps if the NL
$\rightarrow$ formal mapping is unstable. When that mapping is unstable:
\begin{enumerate}
    \item the verification layer has \textbf{no anchor}---it is asked to check an
          object that was never cleanly produced; and
    \item \textbf{hallucinations persist}---an invalid step that survives reduction is
          invisible to a checker that only sees the (malformed) object.
\end{enumerate}
This is consistent with reported experience in neural theorem proving, where the
dominant obstacle is turning an informal or LLM-generated sketch into a form a prover
accepts~\cite{jiang2023draft,yang2023leandojo}, rather than the prover's search power
per se. It also explains why even strong LLMs (e.g., GPT-4) fare poorly as direct
tactic generators on formal libraries~\cite{yang2023leandojo}, while a retrieval-augmented
prover that \emph{anchors} each step in accessible library premises performs far better.

\paragraph{Design implication.}
Investment should be rebalanced toward the reduction layer: better type reconstruction,
disambiguation, and invertibility-preserving encodings. We develop such a design in the
next section.

% ---------------------------------------------------------------------
\section{Proposed Framework}
\label{sec:framework}

\subsection{Typed Intermediate Representation (IR)}
\label{sec:ir}
We propose a single \textbf{Typed IR} as the lingua franca between generation and
verification. It is specified by four components:
\begin{itemize}
    \item \textbf{Type system}: every node carries a type, so ill-formed objects are
          rejected at construction time;
    \item \textbf{Operational semantics}: each construct has a defined meaning, enabling
          deterministic evaluation and symbolic simplification;
    \item \textbf{Constraints}: domain side-conditions (e.g., denominators non-zero,
          domains of definition) are first-class edges in the graph;
    \item \textbf{Compositional structure}: a graph / AST that can be built, split, and
          recomposed modularly.
\end{itemize}
Candidate concrete forms are a DSL, an AST, explicit proof objects, or a workflow
graph; the choice is orthogonal to the four components above.

\subsection{The Pipeline}
\label{sec:pipeline}
Figure~\ref{fig:pipeline} shows the end-to-end pipeline. The key design principle is
that \textbf{natural language is only an interface; the IR is the execution core}.

\begin{figure}[t]
\centering
\fbox{Natural Language (LLM)}
\;\xrightarrow{\text{reduce}}\;
\fbox{Typed IR}
\;\xrightarrow{\text{verify}}\;
\fbox{Formal Verification}
\;\xrightarrow{\text{render}}\;
\fbox{NL Rendering}
\caption{End-to-end pipeline. The typed IR is the execution core that anchors
verification and enables traceability back to the original NL claim.}
\label{fig:pipeline}
\end{figure}

\subsection{Bidirectional Consistency}
\label{sec:bidir}
For the IR to be trustworthy, two properties must hold between IR and NL:
\begin{itemize}
    \item \textbf{Interpretability}: a human can read the IR and understand the reasoning;
    \item \textbf{Traceability}: every IR node maps back to a span of the original NL,
          and every NL claim maps to IR nodes.
\end{itemize}
Bidirectional consistency is what prevents the reduction layer from becoming an
opaque, unaccountable translator---the very failure mode identified in
Section~\ref{sec:bottleneck}.

% ---------------------------------------------------------------------
\section{Related Work}
\label{sec:related}

We situate our framework among four lines of research.

\paragraph{LLM + tools / reasoning strategies.}
Chain-of-thought~\cite{wei2022chain} and tree-of-thoughts~\cite{yao2023tree} improve
generation; Toolformer~\cite{schick2023toolformer} and program-of-thoughts
prompting~\cite{chen2022program} let models call calculators and code interpreters.
These reduce but do not eliminate the verification gap.

\paragraph{Proof assistants and neural theorem proving.}
Early work (GPT-f) trained models to produce proof steps~\cite{polu2020generating};
Hypertree Proof Search scaled neural proving in Lean~\cite{lample2022hypertree};
Draft--Sketch--Prove maps informal proofs to formal sketches that guide a
prover~\cite{jiang2023draft}; LeanDojo/ReProver adds retrieval over the math library
and open tooling~\cite{yang2023leandojo}; autoformalization translates NL statements into
formal ones directly~\cite{wu2022autoformalization}. Our framework generalizes the
\emph{translator} role in these systems into a first-class, typed component.

\paragraph{Neuro-symbolic systems.}
AlphaGeometry combines a neural solver with a symbolic deduction engine for olympiad
geometry~\cite{trinh2024alphageometry}---a existence proof that a neuro-symbolic split,
with the symbolic side as the verifier, can reach human-competitive results.

\paragraph{Learned verification / reward models.}
Outcome verifiers~\cite{cobbe2021training}, process reward models~\cite{lightman2024lets},
and automatic process supervision (Math-Shepherd)~\cite{wang2024mathshepherd} score
intermediate steps. These are verifiers in our sense (Layer~3), and they are exactly why
a stable Layer~2 matters: a process reward model can only supervise steps that have a
well-defined form.

\paragraph{Our difference.}
Prior work treats the NL$\rightarrow$formal translation as a subroutine. We elevate the
\emph{semantic reduction layer} to the central architectural component and argue it is
the dominant bottleneck.

% ---------------------------------------------------------------------
\section{Applications}
\label{sec:apps}

\paragraph{Automated theorem proving (ATP).}
An LLM drafts a proof \emph{sketch}; the sketch is reduced to a typed IR; the IR is
checked in Coq/Lean. This is the pattern of Draft--Sketch--Prove~\cite{jiang2023draft},
which raised success on competition problems from 20.9\% to 39.3\% by anchoring search
in a formal sketch.

\paragraph{Math problem solving.}
A word problem is reduced to a symbolic program (e.g., via program-of-thoughts
prompting~\cite{chen2022program}), executed by a solver, and the result rendered back.
The IR here is the program, which is independently executable and checkable.

\paragraph{AI education.}
A student's answer is reduced to an IR and graded against a reference IR, enabling
transparent, step-level feedback rather than a black-box score.

% ---------------------------------------------------------------------
\section{Evaluation}
\label{sec:eval}

We report \textbf{published baselines} (clearly attributed) to locate the framework, and
we define two metrics that make reliability claims explicit. We do \emph{not} claim a
new large-scale empirical result; the contribution is the framework and the protocol.

\subsection{Benchmarks and Published Baselines}
\label{sec:bench}
Standard benchmarks are GSM8K and the MATH dataset~\cite{hendrycks2021measuring}
for word problems, and MiniF2F~\cite{zheng2021minif2f} / the LeanDojo Benchmark for
formal proving. Representative published numbers:
\begin{itemize}
    \item On GSM8K, process RL with Math-Shepherd lifts Mistral-7B from 77.9\% to
          84.1\%, and verification lifts it further to 89.1\%; on MATH the comparable
          figures are 28.6\% $\rightarrow$ 33.0\% $\rightarrow$ 43.5\%~\cite{wang2024mathshepherd}.
    \item On MiniF2F, ReProver (LeanDojo) achieves 26.5\% Pass@1 and beats GPT-4 as a
          tactic generator; on the LeanDojo Benchmark it reaches 51.2\% (random split)
          and 26.3\% (novel-premise split)~\cite{yang2023leandojo}.
    \item On formal competition problems, Draft--Sketch--Prove improves from 20.9\% to
          39.3\% by using an informal$\rightarrow$formal sketch~\cite{jiang2023draft}.
\end{itemize}
These gains are all \emph{attributable to strengthening a link between generation and
verification}---precisely the reduction layer we foreground.

\subsection{Reliability Metrics Beyond Accuracy}
\label{sec:metrics}
Let a system be evaluated on $N$ problems. Define:
\begin{itemize}
    \item \textbf{Accuracy}
          $A = \frac{1}{N}\sum_{i=1}^{N}\mathbb{1}[\text{final answer of }i\text{ is correct}].$
    \item \textbf{Verifiability}
          $V = \frac{1}{N}\sum_{i=1}^{N}\mathbb{1}[\text{a machine-checkable derivation for }i\text{ exists}].$
    \item \textbf{Consistency} (over the $N_V$ derivable problems)
          $C = \frac{1}{N_V}\sum_{i:\text{ derivable}}\mathbb{1}[\text{NL explanation and formal derivation agree step-wise}].$
\end{itemize}
Accuracy alone can be high while $V$ and $C$ are low (a confident but uncheckable
answer). The framework's promise is to raise $V$ and $C$, bounding what a correctness
claim may assert: a claim is only as reliable as the weakest of $(A, V, C)$.

\begin{table}[t]
\centering
\caption{Reliability metrics and what each bounds.}
\label{tab:metrics}
\begin{tabular}{@{}lll@{}}
\toprule
Metric & Measures & Bound it imposes \\
\midrule
Accuracy ($A$) & final answer correctness & ``Is the output right?'' \\
Verifiability ($V$) & existence of a checkable derivation & ``Can we even check it?'' \\
Consistency ($C$) & NL $\leftrightarrow$ IR agreement & ``Does the story match the proof?'' \\
\bottomrule
\end{tabular}
\end{table}

\subsection{Illustrative Validation}
\label{sec:illustrative}
To show how the reduction layer catches what NL self-checks miss, consider a constructed
example (a \emph{worked illustration}, not a benchmark result). A fluent NL proof might
assert the identity
\begin{equation}
\frac{1}{a-b} = \frac{1}{a} - \frac{1}{b}\qquad(\text{false, in general}).
\end{equation}
A natural-language self-check can be talked into accepting it. Reduced to the typed IR
(an expression tree with the constraint $a\neq b,\ a\neq 0,\ b\neq 0$), a symbolic
simplifier cross-multiplies and obtains $ab = -(a-b)^2$, which is not an identity. The
verification layer therefore \emph{refutes} the claim, whereas the NL layer did not.
This is exactly the failure mode the framework is designed to intercept.

\subsection{A Practical Reliability Checklist}
\label{sec:checklist}
\begin{table}[t]
\centering
\caption{Reliability checklist: threat $\rightarrow$ layer $\rightarrow$ mechanism $\rightarrow$ tool.}
\label{tab:checklist}
\begin{tabular}{@{}llll@{}}
\toprule
Threat & Layer & Mechanism & Example tool \\
\midrule
Wrong final answer & Generation & sampling + search & ToT~\cite{yao2023tree} \\
Untyped / ill-formed & Reduction & type system & Typed IR \\
Ambiguous reading & Reduction & disambiguation & autoformalization~\cite{wu2022autoformalization} \\
Untraceable claim & Reduction & bidirectional consistency & IR $\leftrightarrow$ NL map \\
Invalid step & Verification & proof checking & Lean / Coq \\
Wrong fact & Verification & factual grounding & math library lookup \\
Ungeneralizable & Evaluation & novel-split benchmark & LeanDojo Bench.~\cite{yang2023leandojo} \\
\bottomrule
\end{tabular}
\end{table}

% ---------------------------------------------------------------------
\section{Discussion}
\label{sec:discussion}

\paragraph{Limitations.}
The Typed IR is non-trivial to design, and semantic mapping from NL remains costly;
non-invertibility means perfect traceability is aspirational. We also acknowledge the
framework is a conceptual contribution: we report no new large-scale SOTA and rely on
attributed published numbers.

\paragraph{Scope and boundaries.}
Our claims are bounded. We do \emph{not} assert that every AI4Math task needs full
formal verification---only that any reliability claim should make its $(A,V,C)$ profile
explicit. For low-stakes arithmetic, a calculator call may suffice; for theorem proving,
full formal checking is warranted.

\paragraph{Extensions.}
The three-layer decomposition generalizes beyond mathematics to any domain where NL
must become verified action (code generation, agent planning). Retrieval-augmented
provers~\cite{yang2023leandojo} already hint at agent-style architectures where the IR
node is a tool call.

% ---------------------------------------------------------------------
\section{Conclusion}
\label{sec:conclusion}

We argued that reliable AI reasoning equals \textbf{generation + semantic reduction +
verification}, and that the \emph{semantic reduction layer} is the under-served
bottleneck. We proposed a Typed IR and a four-stage pipeline with bidirectional
consistency, defined verifiability and consistency as explicit reliability metrics, and
distilled a practical checklist. Real systems already validate the pieces; what is
missing is treating reduction as a first-class, typed component. Future work is to build
and empirically stress-test a reference Typed IR across the three application domains.

% ---------------------------------------------------------------------
\bibliographystyle{plain}
\bibliography{references}

\end{document}
